% META: {"id": "1NSI-RES-EVAL-A-corrige", "chapitre": "1NSI-RESEAUX", "type_objet": "corrige_evaluation", "evaluation_ref": "1NSI-RES-EVAL-A", "capacites": ["P-ARCH-02A", "P-ARCH-02B", "P-ARCH-02C", "P-ARCH-04A", "P-ARCH-04B"], "mode_creation": "ex_nihilo", "sources_inspiration": [], "status": "verified"}
% BEGIN-VERIFY
% def decouper_en_paquets(message, taille_max):
%     nb_paquets = (len(message) + taille_max - 1) // taille_max
%     paquets = []
%     for i in range(nb_paquets):
%         donnees = message[i * taille_max:(i + 1) * taille_max]
%         paquets.append({"numero": i, "total": nb_paquets, "donnees": donnees})
%     return paquets
% paquets = decouper_en_paquets("DONNEESNSI", 4)
% assert len(paquets) == 3
% reseau = {"A": ["B"], "B": ["A", "C"], "C": ["B"]}
% def peuvent_communiquer(reseau, depart, arrivee, visites=None):
%     if visites is None:
%         visites = set()
%     if depart == arrivee:
%         return True
%     visites.add(depart)
%     for v in reseau.get(depart, []):
%         if v not in visites and peuvent_communiquer(reseau, v, arrivee, visites):
%             return True
%     return False
% assert peuvent_communiquer(reseau, "A", "C") is True
% END-VERIFY
\begin{corrige}{1NSI-RES-EVAL-A-EX1}
\textbf{1.} $3$ paquets sont créés : \lstinline{"DONN"} (numéro 0), \lstinline{"EESN"}
(numéro 1), \lstinline{"SI"} (numéro 2).

\textbf{2.} Tentative $1$ : envoi de \lstinline{"P"} avec bit $0$ — échoue. Tentative $2$ :
retransmission de \lstinline{"P"} avec bit $0$ — réussit. Tentative $3$ : envoi de
\lstinline{"Q"} avec bit $1$ — réussit. Messages reçus : \lstinline{[(0, "P"), (1, "Q")]}.

\textbf{3.} L'émetteur \textbf{retransmet} systématiquement tant qu'il ne reçoit pas
l'accusé de réception attendu : un message n'est jamais considéré comme envoyé tant que sa
réception n'a pas été confirmée, ce qui garantit qu'il finit par arriver (en l'absence de
perte permanente).
\end{corrige}

\begin{corrige}{1NSI-RES-EVAL-A-EX2}
\textbf{1.} \lstinline{peuvent_communiquer(reseau, "A", "C")} renvoie \lstinline{True} :
il existe un chemin \lstinline{A -> B -> C}.

\textbf{2.}
\begin{python}
def decider_arrosage(humidite, seuil=25):
    return humidite < seuil
\end{python}

\textbf{3.} \lstinline{decider_arrosage(20)} vaut \lstinline{True} (sol sec, il faut
arroser) ; \lstinline{decider_arrosage(30)} vaut \lstinline{False} (sol suffisamment
humide). Le \textbf{capteur} est le capteur d'humidité du sol ; l'\textbf{actionneur} est
la pompe d'arrosage.
\end{corrige}

\begin{corrige}{1NSI-RES-EVAL-A-EX3}
\textbf{1.}
\begin{python}
etat = {"volet_bas": False}

def decider_volet(luminosite, seuil_descente=800, seuil_montee=500):
    if etat["volet_bas"]:
        if luminosite < seuil_montee:
            etat["volet_bas"] = False
    else:
        if luminosite > seuil_descente:
            etat["volet_bas"] = True
    return etat["volet_bas"]
\end{python}

\textbf{2.} \lstinline{[False, True, True, False, False]}. À $700$ lux le volet est déjà
descendu (mesure précédente : $900 > 800$) et $700$ n'est pas \emph{sous} $500$ : il
reste baissé. La décision dépend de l'état mémorisé, pas seulement de la mesure.

\textbf{3.} Non : le cahier des charges dit « dépasse $800$ », donc le test est strict
(\lstinline{luminosite > 800}) et $800$ ne le vérifie pas. Avec un seul seuil, chaque
passage de $790$ à $810$ ferait descendre le volet et chaque retour à $790$ le ferait
remonter : il s'agiterait sans cesse. Les deux seuils ($800$ pour descendre, $500$ pour
remonter) évitent ces basculements autour d'une valeur unique.
\end{corrige}
